
\begin{center}
{\huge \textbf{Escape boat}}\\ \vspace{0.5em}

\renewcommand{\arraystretch}{1.40}
\begin{tabular}{|C{5em}|C{25em}|}
\hline
Cadre  & Intérieur et extérieur du bateau en navigation. \\
\hline
Règles  & Résoudre les différentes étappes le plus vite possible et être le premier bateau à hisser son pavillon ! \\
\hline
Equipes & Équipage de navigation. \\
\hline
Imaginaire & Les explorateurs du ciels zooligistes sont montés sur leur dirigeable à voile afin de partir découvrir de nouvelles espèces. Ils découvrent ce support qui va abriter leur expédition sur les 7 prochains jours. \\
\hline
Rôles & 
Un.e chef.fe met en place le jeu à un moment où les jeunes sont absents (içi pendant la messe).
Le CF lance le jeu au moment opportun. Chef.fe présent.e à bord\\
\hline
Action & Les différentes étappes incluent~:  localisation d'éléments techniques dans le mille piste mer, relèvement de la balise \textit{Basse de chimère}, localisation d'un message chiffré zippé à l'intérieur la boué couronne, récupération d'une clé de déchiffrement qui a été scotchée sur une pinoche, hisser pavillon de couleur. A la fin du jeu les jeunes récupèrent le pavillon de couleur de leur bateau qui a été caché dans la baille de mouillage et le hisse. 
\\
\hline
Sens & Le but est que les jeunes s'approrient le bateau, ses principaux éléments et le matériel de sécurité. \\
\hline
Sécurité & Pas de risques particuliers identifiés à priori. Au moment d'accéder à la baille de mouillage on s'assurera que la personne qui aille sur le pont soit à l'aise.\\
\hline
\end{tabular}
\end{center}

\vspace{1cm}
